% !TEX program = xelatex
\documentclass{article}

% -------------------- Packages --------------------
\usepackage{amsthm, amssymb, bbm, amsmath, mathtools}
\usepackage{graphicx}
\usepackage{xcolor}
\usepackage{booktabs}
\usepackage{multirow}
\usepackage{enumitem}
\usepackage{hyperref}
\usepackage[a4paper, width=5.8in]{geometry}
\usepackage{fontspec}

% -------------------- Fonts --------------------
% \setmainfont{Verdana}
\newcommand{\lmr}{\fontfamily{lmr}\selectfont} % Latin Modern Roman (math-friendly)

% -------------------- Math Macros --------------------
\DeclareMathOperator*{\argmax}{argmax}
\DeclareMathOperator*{\argmin}{argmin}
\newcommand{\R}{\mathbb{R}}
\newcommand{\I}[1]{\mathbbm{1}_{#1}}

% -------------------- Hyperref --------------------
\hypersetup{
    colorlinks=true,
    linkcolor=black,
    urlcolor=blue,
    citecolor=black
}

% -------------------- Title --------------------
\title{Sales Forecasting at Scale: A Machine Learning Study on the Walmart M5 Dataset}
\author{Siddharth Kulkarni}
\date{}

\allowdisplaybreaks

\begin{document}
\maketitle

% -------------------- Abstract --------------------
\begin{abstract}
Accurate demand forecasting is central to inventory management and supply chain efficiency.
This paper presents a machine learning–based sales forecasting pipeline built on the Walmart M5 dataset.
We focus on feature engineering, model selection, and evaluation under realistic business constraints.
Results demonstrate the trade-offs between accuracy, interpretability, and operational complexity.
\end{abstract}

% -------------------- 1. Introduction --------------------
\section{Introduction}
Demand forecasting at scale involves high-dimensional time series, strong seasonality, and sparse item-level data.
The Walmart M5 competition provides a realistic benchmark with hierarchical structure across items, stores, and states \cite{makridakis2022m5}.
This project aims to build an end-to-end forecasting system that balances predictive performance with practical usability.

\textbf{Contributions:}
\begin{itemize}[leftmargin=*]
    \item A structured feature engineering pipeline for large-scale retail time series.
    \item Comparative analysis of classical and ML-based forecasting models.
    \item Business-oriented evaluation using competition-aligned metrics.
\end{itemize}

% -------------------- 2. Dataset --------------------
\section{Dataset Description}
The Walmart M5 dataset contains daily unit sales for over 30,000 products across multiple stores and states.
Auxiliary data includes calendar features, pricing information, and event indicators.

\subsection{Key Characteristics}
\begin{itemize}[leftmargin=*]
    \item Hierarchical structure: item $\rightarrow$ department $\rightarrow$ store $\rightarrow$ state.
    \item Strong weekly and yearly seasonality.
    \item Intermittent demand for many items.
\end{itemize}

% -------------------- 3. Problem Formulation --------------------
\section{Problem Formulation}
Let $y_{i,t}$ denote sales of item $i$ at time $t$.
Given historical observations and covariates $x_{i,t}$, the objective is to learn a forecasting function
\[
\hat{y}_{i,t+h} = f(y_{i,1:t}, x_{i,1:t+h})
\]
for forecast horizon $h$, minimizing a scale-aware loss aligned with business priorities.

% -------------------- 4. Feature Engineering --------------------
\section{Feature Engineering}
Features are constructed using only information available at prediction time.

\subsection{Time-Based Features}
\begin{itemize}[leftmargin=*]
    \item Lagged sales (7, 14, 28 days)
    \item Rolling means and standard deviations
    \item Day-of-week and seasonal indicators
\end{itemize}

\subsection{External Covariates}
\begin{itemize}[leftmargin=*]
    \item Price changes and promotions
    \item Holidays and special events
\end{itemize}

% -------------------- 5. Models --------------------
\section{Models}
We evaluate multiple forecasting approaches:

\begin{itemize}[leftmargin=*]
    \item Baseline statistical models (moving average, exponential smoothing)
    \item Tree-based models (LightGBM)
    \item Global vs item-level modeling strategies
\end{itemize}

Model selection prioritizes robustness and scalability over marginal accuracy gains.

% -------------------- 6. Evaluation --------------------
\section{Evaluation Methodology}
Performance is evaluated using competition-style metrics (e.g., WRMSSE) and standard error measures.
Validation is conducted via rolling-origin splits to respect temporal ordering.

% -------------------- 7. Results --------------------
\section{Results}
Tree-based global models significantly outperform classical baselines, particularly for items with sparse demand.
However, gains diminish for highly seasonal, high-volume products.

\begin{figure}[h]
    \centering
    % \includegraphics[width=0.95\linewidth]{forecast_example.png}
    \caption{Example forecast vs. actual sales}
\end{figure}

% -------------------- 8. Discussion --------------------
\section{Discussion}
While advanced models improve accuracy, they introduce interpretability and maintenance challenges.
Feature leakage and computational cost are key risks in production settings.

% -------------------- 9. Conclusion --------------------
\section{Conclusion}
This study demonstrates a practical ML pipeline for large-scale retail forecasting.
Future work includes hierarchical reconciliation and probabilistic forecasting for inventory optimization.

% -------------------- References --------------------
\bibliographystyle{plain}
\bibliography{references}

\end{document}
